% HD-S29. GPT-6 Astra Pro, /root/sixfolds; inherited effort not exposed.
% 21 November 2026. A complete scoped rank bound with literal matrix maps.
\subsection{HD-S29: a rank bound for the \(G_{12}\)-symmetric tensor source}
\label{sec:g12-tensor-bound}

Put \(K=\mathbb Q(s)\), \(s^2=-2\), and let \(H\subset
\operatorname{GL}_2(K)\) be the group generated by
\[
 R_1=\begin{pmatrix}1&0\\-s&-1\end{pmatrix},\quad
 R_2=\begin{pmatrix}-1&s\\0&1\end{pmatrix},\quad
 R_3=\begin{pmatrix}-1-s/2&-1+s/2\\-1+s/2&1+s/2\end{pmatrix}.
\]
For \(N\ge4\), let \(V_N=\{(x_1,\ldots,x_N)\in K^N:
\sum x_j=0\}\), with the literal coordinate-permutation action of
\(\mathfrak S_N\), and use
\[
 G=H\times\mathfrak S_N,\qquad
 \rho=K^2\otimes_K V_N,\qquad m=\dim_K\rho=2(N-1).
\]
Every connected \(G\)-cover of \(\mathbb P^1\) whose primitive
\(\rho\)-factor is a Weil sixfold has at most nine branch points.
Its period differential therefore has rank at most six. This includes
every specified \(K\)-linear isogeny transport of that factor.
The assertion concerns this actual external tensor representation; it
does not assert a classification of other tensor, induced, or
bireflection representations.

\paragraph{The group and its exact coefficient representation.}
The finite certificate starts with the three displayed matrices,
enumerates all their products, and finds exactly 48 matrices and eight
conjugacy classes. They preserve the positive Hermitian matrix
\[
 H_0=\begin{pmatrix}1&-s/2\\s/2&1\end{pmatrix},
 \qquad \det H_0=\tfrac12.
\]
Every generator is a reflection of order two. The squared absolute
matrix traces sum to 48 over the group, proving absolute
irreducibility by the character inner product. The traces include
\(s\) and are all in \(K\), so the character field is exactly
\(K\). The complete spectrum data are
\[
\begin{array}{c|rrrrrrrr}
\text{class size}&1&1&12&8&6&8&6&6\\
\text{order}&1&2&2&3&4&6&8&8\\
\text{eigenvalues}&1,1&-1,-1&1,-1&
 \zeta_3,\zeta_3^{-1}&i,-i&\zeta_6,\zeta_6^{-1}&
 \zeta_8,\zeta_8^3&\zeta_8^5,\zeta_8^7
\end{array}
\]
with the last two columns exchanged under the other embedding of K.
In particular, every nonidentity element other than a reflection has
both eigenvalues primitive of its own order.

For source identification there is an explicit conjugation. Put
\(r=\sqrt2\), \(s=ir\), and
\[
 T=\begin{pmatrix}1&i\\-1-r&i(1+r)\end{pmatrix}.
\]
The three matrices of Arreche--Williams, \emph{Normal reflection
subgroups of complex reflection groups}, Section 7.3, namely
\[
 r^{-1}\begin{pmatrix}1&-1\\-1&-1\end{pmatrix},\quad
 r^{-1}\begin{pmatrix}1&1\\1&-1\end{pmatrix},\quad
 \begin{pmatrix}0&\zeta_8\\\zeta_8^{-1}&0\end{pmatrix},
\]
are carried to \(R_1,R_2,R_3\) by \(M\mapsto T^{-1}MT\).
This is checked as an equality over \(\mathbb Q(\zeta_8)\).
Thus the representation being tested is literally the published dual
reflection representation of \(G_{12}\), expressed over its
quadratic field; no field-of-definition inference is needed.
The source is DOI \texttt{10.1017/S1474748021000323}, Section 7.3.
Only its displayed candidate matrices are needed: their group order,
spectra, field, and positive form have been recomputed above.

The standard module \(V_N\) is absolutely irreducible: applying
\(1-(ab)\) to any nonscalar coordinate vector produces a nonzero
multiple of \(e_a-e_b\), and its permutations span \(V_N\).
The matrices \(1,R_1,R_2,R_1R_2\) are linearly independent, so
the H-action generates \(M_2(K)\). For the standard action put
\(v_a=e_a-e_N\) and \(\varphi_b(x)=x_b-x_N\). The operator
\(1-(aN)\) equals \(v_a\varphi_a\), and products of these
operators give every \(v_a\varphi_b\), up to the nonzero
coefficient \(\varphi_a(v_b)\in\{1,2\}\). Both the vectors
and the functionals are bases. Hence this action generates
\(\operatorname{End}(V_N)\). The two full matrix algebras tensor
to the full matrix algebra on \(\rho\), proving its absolute
irreducibility. It is faithful. Indeed an element
acting trivially on a tensor product has scalar actions on both factors;
the only scalar permutation on \(V_N\), for \(N\ge3\), is the
identity: it must preserve every unordered pair \(\{a,b\}\)
because it scales \(e_a-e_b\), and all these pairs determine the
individual coordinates. The displayed representation of H is itself faithful.

\paragraph{The exact tensor fixed-space formula.}
Let \(g=(h,\sigma)\) be a local monodromy, let the cycle lengths
of \(\sigma\) be \(a_1,\ldots,a_\ell\), and put
\(f_h(a)=\dim\ker(h^a-1)\). On a cycle of length a, a fixed
vector of \(h\otimes\sigma\) is determined by its first
coefficient x and the equation \(h^a x=x\). The permutation
module splits by the actual projection
\[
 K^N=K(1,\ldots,1)\oplus V_N,\qquad
 (x_j)\longmapsto N^{-1}\Bigl(\sum x_j\Bigr)(1,\ldots,1).
\]
Subtracting the first summand therefore gives the exact formula
\begin{equation}\label{eq:S29-codimension}
 c(h,\sigma):=\operatorname{rank}(1-\rho(h,\sigma))
 =2(N-1)-\sum_{j=1}^{\ell}f_h(a_j)+f_h(1).
\end{equation}
For a reflection h, \(f_h(a)=1+\mathbf1_{2\mid a}\).
Writing e for the number of even cycles, we obtain
\(c=2N-1-\ell-e\ge N-1\), because every odd cycle has
length at least one and every even cycle has length at least two.
For a nonreflection of order o, the spectrum table gives
\(f_h(a)=2\mathbf1_{o\mid a}\) and \(f_h(1)=0\).
Consequently
\[
 c(h,\sigma)\ge2N(1-1/o)-2\ge N-2\qquad(h\ne1).
\]
The reflection case also satisfies the first inequality, and in fact
improves its right side by one. When \(h=1\), the formula is
\(c=2(N-\ell)\); a nonidentity branch element then has
\(\sigma\ne1\) and \(c\ge2\).

\paragraph{The actual H-quotient cannot have genus zero.}
Let \(C\to\mathbb P^1\) be the given connected G-cover and
\(Y=C/\mathfrak S_N\). Its deck group is literally H, and the
projection of a generating branch tuple generates H. Omit projected
identities and let \(o_1,\ldots,o_k\) be the remaining orders.
The group is noncyclic, so \(k\ge3\). Riemann--Hurwitz is
\[
 2g(Y)-2=48\left(-2+\sum_{j=1}^k(1-1/o_j)\right).
\]
Genus zero would give \(\sum(1-1/o_j)=47/24\). Four or
more terms already have sum at least two. For three terms, order them
\(a\le b\le c\). If \(a\ge3\), or if \(a=2,b\ge4\),
the sum is again at least two. For \((a,b)=(2,2)\), the equation
forces \(c=24\), which is absent from the actual order set.
For \((a,b)=(2,3)\), it forces \(c=24/5\), impossible.
Thus \(g(Y)\ge1\), and
\begin{equation}\label{eq:S29-RH}
                  \sum_{j=1}^k(1-1/o_j)\ge2.
\end{equation}

\paragraph{The branch bound and the exact period map.}
For the irreducible representation \(\rho\), the parabolic
cocycle variables are \(c_j\in\operatorname{im}(1-g_j)\),
with the single vector relation \(\sum P_jc_j=0\), modulo
the m-dimensional coboundary space. The relation map is onto: its
annihilator under the positive invariant form would be fixed by every
branch generator, by induction on the prefixes \(P_j\). There
are no invariant vectors. Hence
\[
             \dim_K H^1_{\rm par}=\sum c_j-2m.
\]
For the primitive sixfold this equals six, giving
\(\sum c_j=4N+2\). Here is the actual factor map underlying
that equality. Choose a nonzero vector v in \(\rho(K)\) and
the invariant tensor Hermitian form h. The rational group-algebra
idempotent
\[
 e=\frac{m}{|G|h(v,v)}\sum_{g\in G}
       \operatorname{Tr}_{K/\mathbb Q}h(v,\rho(g^{-1})v)\,g
\]
has rank one on \(\rho\), rank one on its conjugate, and zero
on the other complex irreducible blocks, by matrix-coefficient
orthogonality. Clearing its denominators gives an actual endomorphism
\(De\) of the Jacobian using covariant deck maps. Its connected
image B is the primitive factor. The coefficient maps into cover
cohomology identify its first cohomology with the parabolic
multiplicity space and its conjugate, as in the explicit averaging
construction of HD-S23. Thus \(\dim_K H^1(B)=6\) is exactly
the above parabolic dimension. This is a morphism of the fixed relative
Jacobians over the marked Hurwitz family.

Let b be the total number of true branch points and k the number with
nontrivial H projection. For \(k\ge4\), the inequalities above
give, since \(N\ge4\),
\[
 4N+2=\sum c_j\ge k(N-2)+2(b-k)
       =k(N-4)+2b\ge4(N-4)+2b.
\]
Therefore \(b\le9\). For \(k\le3\), combine
\eqref{eq:S29-RH} with the order-dependent bound to obtain
\[
 4N+2\ge2N\sum_{j=1}^k(1-1/o_j)-2k+2(b-k)
       \ge4N+2b-4k,
\]
so \(b\le2k+1\le7\). These two cases prove the asserted
universal bound \(b\le9\).

Finally the marked Hurwitz family maps finitely and locally to the
configuration of b distinct ordered branch points: its discrete fibre
is the finite set of generating product-one tuples, and retaining the
sheet permutations while moving branch cuts gives the inverse local
construction. Quotient by the three-dimensional target automorphism
group gives dimension \(b-3\le6\). The period map of the
actual relative image of De factors through this family, proving the
differential-rank bound. A specified K-linear isogeny is a fixed
invertible rational matrix on local homology markings; its map on
Hodge subspaces has the inverse induced by the inverse matrix.
It preserves this rank bound.

\paragraph{Check record.}
The exact GAP file \path{g12_matrix_certificate.g} reconstructs the
literal matrices, published conjugator, 48-element closure, eight
classes, positive form and every spectrum. The Python file
\path{g12_tensor_bounds.py} independently evaluates the resulting
fixed-space formula on every permutation cycle partition for
\(4\le N\le20\), and checks the finite genus-zero equation.
Those finite checks supplement the uniform inequalities proved above;
they do not replace them. The complete author proof has been read at
the end of writing. Independent review is recorded separately when it
occurs.
